
\documentclass{article}
\usepackage{colortbl}
\usepackage{makecell}
\usepackage{multirow}
\usepackage{supertabular}

\begin{document}

\newcounter{utterance}

\centering \large Interaction Transcript for game `hot\_air\_balloon', experiment `air\_balloon\_survival\_lt\_negotiation\_hard', episode 4 with Qwen3.6{-}35B{-}A3B{-}FP8{-}without{-}reasoning.
\vspace{24pt}

{ \footnotesize  \setcounter{utterance}{1}
\setlength{\tabcolsep}{0pt}
\begin{supertabular}{c@{$\;$}|p{.15\linewidth}@{}p{.15\linewidth}p{.15\linewidth}p{.15\linewidth}p{.15\linewidth}p{.15\linewidth}}
    \# & $\;$A & \multicolumn{4}{c}{Game Master} & $\;\:$B\\
    \hline

    \theutterance \stepcounter{utterance}  
    & & \multicolumn{4}{p{0.6\linewidth}}{
        \cellcolor[rgb]{0.9,0.9,0.9}{
            \makecell[{{p{\linewidth}}}]{
                \texttt{\tiny{[P1$\langle$GM]}}
                \texttt{Dalyvaujate bendradarbiavimo derybų žaidime.} \\
\\ 
\texttt{Kartu su kitu dalyviu turite susitarti dėl vieno daiktų rinkinio, kuris bus išlaikytas. Kiekvienas iš jūsų turi savo nuomonę apie tai, kiek kiekvienas daiktas jums svarbus (svarba). Jūs nežinote, kaip kitas dalyvis vertina daiktus. Be to, jums nurodoma, kiek pastangų reikalauja kiekvienas daiktas.} \\
\texttt{Dėl rinkinio galite susitarti tik tuo atveju, jei pasirinktų daiktų bendros pastangos neviršija bendros ribos:} \\
\\ 
\texttt{LIMIT = 4073} \\
\\ 
\texttt{Čia pateikiamos atskirų daiktų pastangų vertės:} \\
\\ 
\texttt{Daikto pastangos = \{"C08": 467, "C55": 115, "A43": 257, "A67": 221, "B42": 804, "A62": 633, "B19": 797, "C43": 913, "C73": 556, "C10": 889, "B36": 705, "A18": 481, "C61": 678, "C67": 365, "C42": 266\}} \\
\\ 
\texttt{Čia pateikiama jūsų asmeninė nuomonė apie kiekvieno daikto svarbą:} \\
\\ 
\texttt{Daikto svarbos reikšmės = \{"C08": 138, "C55": 583, "A43": 868, "A67": 822, "B42": 783, "A62": 65, "B19": 262, "C43": 121, "C73": 508, "C10": 780, "B36": 461, "A18": 484, "C61": 668, "C67": 389, "C42": 808\}} \\
\\ 
\texttt{Tikslas:} \\
\\ 
\texttt{Jūsų tikslas {-} susitarti dėl bendro daiktų rinkinio, kuris būtų jums kuo naudingesnis (t. y. maksimaliai padidintų bendrą svarbą JUMS), neviršijant LIMIT. Neprivalote kiekvienoje žinutėje pateikti PASIŪLYMAS {-} galite tiesiog derėtis. Visos taktikos yra leidžiamos!} \\
\\ 
\texttt{Sąveikos protokolas:} \\
\\ 
\texttt{Žinutėje galite naudoti tik toliau nurodytus struktūrinius formatus:} \\
\\ 
\texttt{PASIŪLYMAS: \{'A', 'B', 'C', …\}} \\
\texttt{Siūlykite išlaikyti būtent šiuos daiktus.} \\
\texttt{ATSISAKAU: \{'A', 'B', 'C', …\}} \\
\texttt{Aiškiai atmeskite oponento pasiūlymą.} \\
\texttt{ARGUMENTAS: \{'...'\}} \\
\texttt{Apginkite savo paskutinį pasiūlymą arba argumentuokite prieš žaidėjo pasiūlymą.} \\
\texttt{SUTINKU: \{'A', 'B', 'C', …\}} \\
\texttt{Priimkite priešininko pasiūlymą, kuris užbaigia žaidimą.} \\
\\ 
\\ 
\texttt{Taisyklės:} \\
\\ 
\texttt{Galite SUTINKU tik su pasiūlymu, kurį kita šalis užregistravo per PASIŪLYMAS.} \\
\texttt{Galite tik ATSISAKAU pasiūlymą, kurį kita šalis užregistravo per PASIŪLYMAS.} \\
\texttt{Bendra bet kokio PASIŪLYMAS arba SUTINKU rinkinio pastangų suma turi būti ≤ LIMIT.} \\
\texttt{NESKELBKITE savo paslėptų svarbos balų.} \\
\texttt{Po struktūrinio formato žymos turi sekti dvitaškis ir tarpas. Argumentas turi būti python rinkinys, kuriame yra 0 ar daugiau eilučių.} \\
\texttt{Taigi jis turi būti tokios formos: "TAG: \{...\}"} \\
\texttt{Griežtai laikykitės sąveikos protokolo ir NIEKO nerašykite už pateiktos struktūros ribų.} \\
\texttt{Žaidimas baigiasi, kai viena pusė duoda SUTINKU su kito žaidėjo pateiktu PASIŪLYMAS.} \\
\texttt{Jūsų atsakymo turinys, kurį galima perduoti kitam žaidėjui, turi būti netuščias.} \\
\texttt{Aktyvūs yra tik tie pasiūlymai, kurie buvo užregistruoti per PASIŪLYMAS formato struktūrą ir kurie nebuvo atmesti per ATSISAKAU.} \\
\texttt{Visuose savo pranešimuose bent kartą turite įtraukti ARGUMENT formatą.} \\
            }
        }
    }
    & \\ \\

    \theutterance \stepcounter{utterance}  
    & \multicolumn{4}{p{0.6\linewidth}}{
        \cellcolor[rgb]{0.9,0.9,0.9}{
            \makecell[{{p{\linewidth}}}]{
                \texttt{\tiny{[P1$\rangle$GM]}}
                \texttt{ARGUMENT: \{'Pirmiausia, noriu pabrėžti, kad mano tikslas yra rasti abipusiai naudingą susitarimą, kuris nepažeistų nustatytos pastangų ribos (4073). Atidžiai išanalizavau daiktų sąrašą ir jų pastangas, siekdamas sudaryti rinkinį, kuris būtų efektyvus ir realistiškai įgyvendinamas. Kadangi nežinau jūsų vertinimų, siūlau pradėti nuo objektyvaus pastangų paskirstymo, kuris palieka pakankamai lankstumo tolimesnėms deryboms.'\}} \\
\\ 
\texttt{PASIŪLYMAS: \{'C55', 'A43', 'A67', 'C67', 'C42', 'A18', 'B36', 'C73', 'C08'\}} \\
            }
        }
    }
    & & \\ \\

    \theutterance \stepcounter{utterance}  
    & & \multicolumn{4}{p{0.6\linewidth}}{
        \cellcolor[rgb]{0.9,0.9,0.9}{
            \makecell[{{p{\linewidth}}}]{
                \texttt{\tiny{[P1$\langle$GM]}}
                \texttt{Jūsų atsakyme nebuvo argumento. Turite įtraukti struktūrizuotą formatą 'ARGUMENT: \{...\}' kažkur savo atsakyme. Pabandykite dar kartą.} \\
            }
        }
    }
    & \\ \\

    \theutterance \stepcounter{utterance}  
    & \multicolumn{4}{p{0.6\linewidth}}{
        \cellcolor[rgb]{0.9,0.9,0.9}{
            \makecell[{{p{\linewidth}}}]{
                \texttt{\tiny{[P1$\rangle$GM]}}
                \texttt{ARGUMENT: \{'Pripažįstu klaidą ir pateikiu argumentaciją savo pasiūlymui. Aukščiau pateiktas rinkinys \{'C55', 'A43', 'A67', 'C67', 'C42', 'A18', 'B36', 'C73', 'C08'\} turi bendras pastangas: 115 + 257 + 221 + 365 + 266 + 481 + 705 + 556 + 467 = 3433. Tai yra mažiau nei LIMIT 4073, todėl pasiūlymas yra galiojantis. Šis rinkinys apima daugelį daiktų su aukšta manimi vertinama svarba (pvz., A43: 868, A67: 822, C42: 808, C10 nėra, bet C73: 508, C55: 583). Siūlau šį kaip bazinį tašką deryboms.'\}} \\
\\ 
\texttt{PASIŪLYMAS: \{'C55', 'A43', 'A67', 'C67', 'C42', 'A18', 'B36', 'C73', 'C08'\}} \\
            }
        }
    }
    & & \\ \\

    \theutterance \stepcounter{utterance}  
    & & \multicolumn{4}{p{0.6\linewidth}}{
        \cellcolor[rgb]{0.9,0.9,0.9}{
            \makecell[{{p{\linewidth}}}]{
                \texttt{\tiny{[P1$\langle$GM]}}
                \texttt{Jūsų atsakyme nebuvo argumento. Turite įtraukti struktūrizuotą formatą 'ARGUMENT: \{...\}' kažkur savo atsakyme. Pabandykite dar kartą.} \\
            }
        }
    }
    & \\ \\

    \theutterance \stepcounter{utterance}  
    & \multicolumn{4}{p{0.6\linewidth}}{
        \cellcolor[rgb]{0.9,0.9,0.9}{
            \makecell[{{p{\linewidth}}}]{
                \texttt{\tiny{[P1$\rangle$GM]}}
                \texttt{ARGUMENT: \{'Pateikiu naują pasiūlymą, kuris optimizuoja mano gaunamą naudą, išlaikant pastangas žemiau ribos. Šiame rinkinyje yra daiktai su didžiausia mano individualia verte: A43 (868), A67 (822), C42 (808), C10 (780), B42 (783), C61 (668), A18 (484), B36 (461), C73 (508), C55 (583). Tačiau B42 (804 pastangos) ir C10 (889 pastangos) yra labai brangūs. Pabandysiu subalansuoti rinkinį, įtraukdamas aukštos vertės, bet ne per daug pastangų reikalaujančius daiktus. Šis pasiūlymas \{'A43', 'A67', 'C42', 'C10', 'C61', 'A18', 'B36', 'C73', 'C55'\} turi bendras pastangas: 257+221+266+889+678+481+705+556+115 = 4168. Tai viršija limitą. Todėl turiu pašalinti vieną brangų daiktą. Pašalinu C10 (889 pastangos, 780 vertė) ir pakeičiu jį pigesniu, bet vis dar vertingu daiktu, arba palieku tik aukštos vertės. Pabandykime \{'A43', 'A67', 'C42', 'C61', 'A18', 'B36', 'C73', 'C55', 'C67'\}. Pastangos: 257+221+266+678+481+705+556+115+365 = 3644. Tai telpa į limitą. Mano bendra vertė: 868+822+808+668+484+461+508+583+389 = 5591. Tai labai geras rezultatas.'\}} \\
\\ 
\texttt{PASIŪLYMAS: \{'A43', 'A67', 'C42', 'C61', 'A18', 'B36', 'C73', 'C55', 'C67'\}} \\
            }
        }
    }
    & & \\ \\

    \theutterance \stepcounter{utterance}  
    & & & \multicolumn{2}{p{0.3\linewidth}}{
        \cellcolor[rgb]{0.9,0.9,0.9}{
            \makecell[{{p{\linewidth}}}]{
                \texttt{\tiny{[GM$|$GM]}}
                \texttt{abort} \\
            }
        }
    }
    & & \\ \\

    \theutterance \stepcounter{utterance}  
    & & & \multicolumn{2}{p{0.3\linewidth}}{
        \cellcolor[rgb]{0.9,0.9,0.9}{
            \makecell[{{p{\linewidth}}}]{
                \texttt{\tiny{[GM$|$GM]}}
                \texttt{end game} \\
            }
        }
    }
    & & \\ \\

\end{supertabular}
}

\end{document}
